\section{从手工特征到自监督：模型究竟在学习什么}

这堂 V6 开场课并不把 Transformer 简化成一条 attention 公式。两位讲者先追溯学习信号如何从手工特征、监督标签发展到自监督目标，再把数据、post-training、agent、multimodality、memory、alignment 与替代架构接到同一张系统地图上。读这份讲义时应持续追问三个问题：信息从哪里来，梯度奖励了什么，部署后的系统还增加了哪些模型外机制？

\subsection{课程承诺：机制、应用与边界必须一起看}

开场的 takeaways 把 attention、scaling、跨域应用、frontier speakers 与 open problems 放在同一页。这个顺序定义了讲义的组织原则：先解释机制，再核对证据，最后写清边界。仅知道一个 block 的张量形状，不足以解释模型为何在某个数据集上成功；仅看到成功 demo，也不足以判断它在新分布、长上下文或真实工作流中是否可靠。

\lecturefigure{slide-011.jpg}{课程目标：从 Transformer 机制到前沿应用、限制与开放问题}{CS25 V6 official deck, p.~11}

\teachervoice{00:06:12--00:07:10，老师强调课程不只讲 attention，而是要让学生接触前沿应用、关键限制和仍可贡献的开放问题。课堂提示：后文每种方法都应回答“改变了系统的哪一层”。}

\begin{importantbox}{现代 Transformer 系统的六层接口}
\begin{enumerate}
  \item \textbf{Representation}：原始输入如何变成 token 与向量；
  \item \textbf{Architecture}：信息如何在序列、层和模态之间流动；
  \item \textbf{Pretraining}：数据分布与目标函数塑造哪些默认能力；
  \item \textbf{Post-training}：监督、偏好、搜索与 verifier 如何修改行为；
  \item \textbf{System loop}：检索、工具、memory 与环境反馈怎样进入闭环；
  \item \textbf{Governance/evidence}：评估、解释、alignment 与部署边界如何定义。
\end{enumerate}
一个新名词若不能落到这些接口之一，通常还没有被真正解释。
\end{importantbox}

\subsection{三次转变：特征工程、监督表示与自监督表示}

早期机器学习依赖研究者手工设计特征，再由一个相对浅的模型把特征映射到标签。深度学习把表示学习交给网络，但仍要求大量人工标签。自监督学习继续前移监督来源：标签不再来自人工类别，而是从输入自身构造，例如遮住 token、预测下一步、恢复被扰动图像或区分匹配与不匹配的样本。三种范式的关键差别不是“网络深不深”，而是\emph{谁定义了学习任务以及模型能看到多少原始结构}。

\lecturefigure{slide-013.jpg}{早期机器学习：手工特征、优化器与昂贵人工标签}{CS25 V6 official deck, p.~13}

\lecturefigure{slide-014.jpg}{监督深度学习：网络开始承担表示学习，但仍依赖标签}{CS25 V6 official deck, p.~14}

从 slide 13 到 slide 16 应按信息瓶颈阅读：手工特征把原始数据压成研究者认为有用的维度，监督网络扩大了可学习表示的范围，而端到端训练进一步减少人为中间接口。代价也随之变化：模型自由度提高，需要更多数据与 compute；如果训练分布带有偏差，网络会把偏差一起学进去。

\lecturefigure{slide-015.jpg}{监督学习的过渡状态：原始数据仍先经过手工表示}{CS25 V6 official deck, p.~15}

\lecturefigure{slide-016.jpg}{端到端监督学习：原始输入直接进入深层网络}{CS25 V6 official deck, p.~16}

\begin{knowledgebox}{读图：为什么四页渐进图都值得保留}
先比较谁在构造 representation，再比较监督信号从哪里来。slide 13 中人类定义特征；slide 14--15 中网络与手工接口共同存在；slide 16 中网络直接消费 raw data。图支持“表示学习接口逐渐内移”，却不证明人工知识没有价值。现代系统仍通过 tokenizer、数据过滤、架构先验、loss 与工具 API 注入大量人为设计。
\end{knowledgebox}

自监督学习的目标可抽象为从样本 $x$ 构造可预测部分 $y=g(x)$ 与可见上下文 $\tilde{x}=c(x)$，然后优化
\[
\mathcal{L}_{\mathrm{SSL}}(\theta)
=-\log p_{\theta}\!\left(y\mid \tilde{x}\right).
\]
$g$ 表示构造目标的规则，例如下一 token 或被遮蔽 patch；$c$ 表示保留给模型的上下文；$\theta$ 是参数。监督虽然“来自数据本身”，但 $g$ 与 $c$ 仍决定模型被鼓励保留什么信息。

\lecturefigure{slide-017.jpg}{自监督学习：从原始数据内部构造训练目标}{CS25 V6 official deck, p.~17}

\teachervoice{00:07:10--00:08:13，讲者把历史主线概括为 hand-engineered features、supervised representation learning 与 self-supervision。老师强调自监督减少人工标签，却没有消除 objective 所带来的 inductive bias。}

\begin{warningbox}{“没有人工标签”不等于“没有人为假设”}
遮蔽比例、预测方向、负样本、数据增强和窗口长度都会改变可学习信号。若增强破坏了任务相关因素，contrastive model 会被奖励去忽略真正有用的信息；若 next-token 数据充满模板与重复，模型会优先学习高频形式。自监督是标签构造方式，不是价值中立的自然规律。
\end{warningbox}

从 information bottleneck 的角度看，objective 决定哪些信息值得保留。情感分类只需要保留与标签相关的方向，可能主动丢弃实体、时序和语法细节；重建目标需要保留更多局部信息，却可能把大量像素噪声也编码进去；next-token prediction 要压缩对未来 token 有帮助的统计规律，但对不会改变 token 概率的真实世界状态没有直接奖励。评估 representation 时因此不能只看一个 downstream score，而应检查线性 probe、few-shot transfer、distribution shift 与可干预性。一个表示在分类上好，不代表它适合规划；一个生成模型能重建细节，也不代表其 latent 已经抽象出对象和因果关系。

\subsection{语言与视觉案例：预测目标决定表示的颗粒度}

语言案例从情感分类转向 next-token prediction。前者只要求把一段话压成少数标签，后者要求在每个位置预测条件分布，因此监督更密集。设 token 序列为 $x_{1:T}$，自回归语言模型优化
\[
\mathcal{L}_{\mathrm{NTP}}(\theta)
=-\sum_{t=1}^{T}\log p_{\theta}(x_t\mid x_{<t}).
\]
$x_t$ 是第 $t$ 个 token，$x_{<t}$ 是历史上下文。每个位置都贡献损失，但目标仍只要求预测观测序列，并不直接要求模型给出事实来源、执行工具或维护长期记忆。

\lecturefigure{slide-018.jpg}{语言案例：从稀疏情感标签转向密集序列预测}{CS25 V6 official deck, p.~18}

\lecturefigure{slide-020.jpg}{语言自监督的完整状态：上下文条件下预测下一个 token}{CS25 V6 official deck, p.~20}

\teachervoice{00:08:13--00:09:57，老师用 sentiment classification 对比 language modeling：分类只提供一个低带宽标签，而序列预测在每个位置产生监督。实践经验：监督更密集不代表目标自动覆盖事实性与因果性。}

视觉自监督通常在 reconstruction 与 contrastive learning 之间选择。重建目标要求 decoder 恢复输入，容易保留像素细节；contrastive objective 则让同一对象的增强视图靠近、不同样本远离，更强调不变表示。一个简化的 InfoNCE loss 是
\[
\mathcal{L}_{\mathrm{NCE}}
=-\log\frac{\exp(\operatorname{sim}(z_i,z_i^+)/\tau)}
{\sum_j \exp(\operatorname{sim}(z_i,z_j)/\tau)}.
\]
$z_i$ 是 anchor 表示，$z_i^+$ 是正样本，$z_j$ 是候选样本，$\tau$ 是温度。关键不是记住公式，而是理解正负样本定义决定了模型被要求忽略哪些变化。

\lecturefigure{slide-021.jpg}{视觉案例：autoencoder 重建与 contrastive learning 的两种监督接口}{CS25 V6 official deck, p.~21}

\lecturefigure{slide-022.jpg}{Masked Autoencoder：遮蔽 patch 后预测缺失视觉内容}{CS25 V6 official deck, p.~22}

\begin{knowledgebox}{读图：reconstruction 与 contrastive 的证据边界}
先看输出空间：重建模型预测输入细节，contrastive model 预测样本关系。再看 invariance：后者会主动压掉增强中被视为“不重要”的变化。图支持两类 representation objective 的差异，却不能单凭示意图断言哪种目标普遍更好；下游任务、数据规模、augmentations 和 decoder 容量都会改变结论。
\end{knowledgebox}

\subsection{本章小结}

学习范式的演化可以压缩为“监督接口不断内移”：先由人类造特征，再由网络学表示，最后由数据自身生成密集目标。但任何目标都只奖励特定统计规律。理解 Transformer 之前，必须先理解它接收的是 token、优化的是条件预测，而非自动获得事实、行动与长期一致性。

